\documentclass[sigplan,10pt]{acmart}

\settopmatter{printacmref=false}
\renewcommand\footnotetextcopyrightpermission[1]{}
\pagestyle{plain}
\settopmatter{printfolios=true}

\usepackage[utf8]{inputenc}
\usepackage[T1]{fontenc}
\usepackage{booktabs}
\usepackage{tabularx}
\usepackage{array}
\usepackage{url}
\usepackage{xspace}
\usepackage{microtype}
\hypersetup{colorlinks=true,linkcolor=black,citecolor=black,urlcolor=black}

\newcommand{\sysname}{OmniFlow\xspace}

\title{Expanded Background and Related Work for \sysname}
\author{}
\date{}

\begin{document}
\maketitle

\section{Background: Mobile GUI Agents as Repeated Observe-Reason-Act Systems}

Recent mobile GUI agents usually execute tasks through an
observe-reason-act loop. At each step, the system observes the current screen
through a screenshot, an accessibility tree, or both; constructs a prompt or
structured state; asks an LLM or VLM to choose the next action; dispatches a
tap, swipe, text input, back action, or app-level command; and repeats after the
environment changes. AndroidWorld formalizes this setting as a dynamic Android
benchmark with programmatic initialization, success checking, and teardown logic
for 116 tasks across 20 real-world Android apps~\cite{rawles2025androidworld}.
Android in the Wild (AITW) provides a much larger demonstration corpus with
715K episodes, around 30K unique instructions, Android versions 10--13, and
eight Pixel device types~\cite{rawles2023android}. These benchmarks make two
facts clear. First, GUI agents need semantic interpretation of natural language
and visual state. Second, the execution substrate is heterogeneous: device size,
OS version, app state, account data, and available widget metadata all affect
whether a recorded action can still be used.

The cost profile of this loop is important for systems work. A step-wise agent
pays model latency and token cost every time it asks the model to select the
next action. V-Droid reports that prior high-performing agents on AndroidWorld
can take tens of seconds per step, while its verifier-based design reaches
4.3\,s per step and 59.5\% task success on
AndroidWorld~\cite{dai2025advancing}. The key point for \sysname is that this is
still per-step online decision making: the agent reduces the cost of each
decision, yet a repeated routine still passes through the model or verifier for
each screen transition. AndroidWorld also reports that its best baseline in the
original benchmark completed 30.6\% of tasks, leaving significant headroom for
robust mobile control~\cite{rawles2025androidworld}. In practice, even a fast
per-step policy creates cumulative cost for long workflows, especially when the
same or similar routines are repeated by many users.

Repeated mobile behavior is not a corner case. Smartphone-use studies describe
brief, repetitive checking behavior as a common usage
pattern~\cite{oulasvirta2012habits}. App launching itself is sufficiently
predictable that FALCON used location and temporal access patterns to prelaunch
likely apps and reduce perceived launch delay~\cite{yan2012fast}. AITW's 715K
episodes over roughly 30K unique instructions further suggests that
device-control workloads contain repeated task forms even when the concrete
screen state and parameters vary~\cite{rawles2023android}. This motivates a
cache-system view of GUI agents. If a user, device, or prior execution has
already produced a verified route through a repeated GUI segment, later
executions should be able to reuse that segment without asking the online model
to reinterpret every screen.

The challenge is that GUI reuse is harder than replaying an API call. Android
accessibility trees expose bounds and hierarchy reliably when elements are
visible, while semantic fields such as \texttt{resource-id}, \texttt{text}, and
\texttt{content-desc} depend on developer annotation, framework behavior, and
app implementation. WebView, custom-rendered views, obfuscated identifiers,
dynamic list content, dialogs, IME overlays, and account-specific state make raw
coordinate replay and exact selector replay fragile. A UI transition graph
(UTG) captures useful structure: nodes represent GUI states and edges represent
actions or transitions. However, a graph built from raw XML snapshots or static
window transitions records possible movement through an interface; it does not
by itself decide whether a historical action has the same task effect in the
current state. For \sysname, the reusable object therefore has to combine state
localization, target transfer, and effect checking. The paper's current
abstraction is a \texttt{Function}: a reusable executable task segment tied to a
start condition, action evidence, transfer logic, and terminal/effect
verification. \texttt{PageVectorSet} provides a privacy-bounded state key for
candidate recall, and \texttt{run\_log -> function asset} turns successful
executions into future cache entries.

\section{Related Work}

\subsection{Record-and-Replay and Programming by Demonstration}

Classic mobile automation systems show the value of executable reuse. RERAN
records low-level Android input events and replays timing-sensitive gestures and
sensor events~\cite{gomez2013reran}. Its strength is fidelity: it captures the
real event stream rather than a symbolic approximation of user input. This makes
it valuable for reproducing executions during testing or debugging. Its
limitation follows from the same design. The replay is tied to the recorded
device state, timing, layout, and app behavior. When an element shifts, a dialog
appears, or a list item moves, replaying the same event stream can target the
wrong location or fail to produce the intended task effect.

SUGILITE takes a higher-level programming-by-demonstration approach for
smartphone automation~\cite{li2017sugilite}. A user demonstrates a task directly
in Android apps, and the system uses the accessibility API to construct an
automation script that can later be invoked through natural language. This line
of work is important because it treats repeated GUI work as something that users
should not have to reperform manually. It also shows a limitation that remains
central for GUI-agent cache systems: demonstrations must be generalized from
observed widgets and states to future states. If the future UI lacks stable
attributes or appears in a different layout/context, the demonstration needs
repair, clarification, or a fallback path.

DroidBot and related GUI-testing tools construct state-transition models during
exploration~\cite{droidbot}. DroidBot can generate a UTG from UI-guided testing
and use that graph to guide future exploration. Static analysis work such as
GATOR's window transition graph estimates possible Android GUI transitions from
application code~\cite{gator}. These systems provide graph structure, but their
objective is testing or static modeling rather than executing user goals from
natural language under dynamic state. Static graphs can over-approximate
infeasible transitions or miss runtime-only transitions; dynamic exploration
graphs depend on coverage and collected app states. \sysname uses a graph
vocabulary because state and transition structure are useful, while the cache
object is a verified \texttt{Function} with start-node preconditions and
terminal/effect checks.

\subsection{Memory-Based Mobile GUI Agents}

AutoDroid is a representative memory-augmented mobile GUI
agent~\cite{wen2024autodroid}. Its method combines LLM commonsense with
app-specific knowledge collected through automated dynamic analysis. The system
builds a functionality-aware UI representation, injects exploration-based
memory to augment the LLM with app knowledge, and uses multi-granularity query
optimization to reduce inference cost. It reports 90.9\% action accuracy and
71.3\% task success on a benchmark of 158 common tasks, outperforming
GPT-4-powered baselines by 36.4\% in action accuracy and 39.7\% in task
success. AutoDroid therefore demonstrates that app memory and structured UI
representation can materially improve task automation.

The limitation for \sysname's target problem is the execution boundary.
AutoDroid still treats memory mainly as knowledge for the LLM's decision
process. The memory helps the model understand app-specific states and actions,
but a repeated task is not necessarily served as a directly executable cache hit
that skips online reasoning for the matched segment. The app-specific
exploration phase also creates a cold-start cost and a deployment question:
users or hosts need a way to collect, store, and update app memory before reuse
becomes available. For a cache-first system, the reusable unit has to be
callable by the execution runtime, guarded by state/effect evidence, and written
back from accepted run logs instead of existing only as prompt context.

MobileGPT makes the reuse goal more explicit through a human-like app memory
pipeline~\cite{lee2024mobilegpt}. It frames mobile task learning as explore,
select, derive, and recall. The agent explores an app, selects relevant
knowledge, derives reusable subtask procedures, and recalls them for later
tasks. It evaluates on 185 tasks across 18 apps and reports 82.7\% accuracy for
automating and learning new tasks, 98.75\% adaptation accuracy across contexts,
and latency/cost reductions of 62.5\% and 68.8\% relative to a GPT-4 baseline.
This is strong evidence that memory over prior task experience can reduce
repeated cost.

MobileGPT's limitation is also instructive. A human-like memory stores
procedures and app knowledge that the LLM can recall and adapt. The online agent
still needs to bind recalled knowledge to the current screen, and the memory is
not the same object as a guarded executable cache entry. When UI attributes are
missing, list content changes, or app state diverges from the learned procedure,
the system has to adapt through the current reasoning loop. \sysname's
\texttt{Function} boundary targets a narrower but more directly executable form
of reuse: if the current state satisfies the start precondition and transfer
checks, the runtime executes the cached segment under
\texttt{control-aware execution}; if checks fail, the host agent remains the
fallback.

AppAgent uses autonomous exploration or human demonstrations to build an app
knowledge base for a multimodal smartphone agent~\cite{zhang2025appagent}. It
operates through a simplified action space that mimics user interactions such as
taps and swipes, and it was tested on 50 tasks across 10 applications. AppAgent
is relevant because it shows that exploration documents and demonstrations can
help a general multimodal agent operate unfamiliar apps. Its reuse unit,
however, remains knowledge for future decisions. It does not define a shared,
privacy-bounded \texttt{PageVectorSet} state key or a
\texttt{run\_log -> function asset} lifecycle in which verified executions
become executable cache assets.

\subsection{Programmatic and State-Machine GUI Agents}

Programmatic GUI agents move beyond prompting by converting task knowledge into
executable code or state-machine plans. AutoDroid-V2 adopts a full task-level
script generation approach for mobile GUI agents~\cite{autodroid-v2}. Its
motivation is to make small-language-model GUI agents more effective by
generating code that can execute task logic. This direction is useful for
repeated workflows because executable scripts can reduce repeated
natural-language reasoning. The tradeoff is that task-level scripts need a
stable enough app model, current-state alignment, and error handling when the
live UI deviates from the generated assumptions. A complete task script is also
a relatively coarse reuse object: if the beginning of a workflow changes while
the effect-carrying middle segment is still valid, the script needs repair or
regeneration.

ActionEngine is a recent state-machine-memory system that shifts GUI agents
from reactive step-by-step operation to programmatic
execution~\cite{zhong2026actionengine}. It uses a Crawling Agent to construct an
updatable state-machine memory through offline exploration and an Execution
Agent to synthesize complete Python programs for online task execution. On
Reddit tasks from WebArena, it reports 95\% task success with about one LLM call
on average, 11.8$\times$ cost reduction, and 2$\times$ end-to-end latency
reduction compared with a strong vision-only baseline. ActionEngine is close to
\sysname in recognizing that a GUI can be reused through a state-machine-like
memory rather than reinterpreted at every step.

The boundary remains different. ActionEngine's reported result is a web task
setting, and the method depends on an offline crawler building a state-machine
memory from the target GUI. Execution failures trigger vision-based re-grounding
and memory updates. \sysname's mobile setting has stronger device heterogeneity,
privacy, and app-state constraints. The paper's current contribution should
therefore avoid claiming a broad superiority over state-machine agents. The
precise gap is that \sysname makes the cache unit a \texttt{Function} derived
from accepted run logs, uses \texttt{PageVectorSet} as the privacy-bounded
state-recall object, and requires start-node, route-only bridge, action
transfer, and terminal/effect verification before a cache hit is counted as
successful.

AgentProg addresses long-horizon GUI tasks from the context-management
angle~\cite{tian2025agentprog}. It represents interaction history as a program
with variables and control flow, and it adds a global belief state to handle
partial observability and environmental change. This is important because long
GUI tasks quickly accumulate history that can overflow or dilute model context.
AgentProg's limitation relative to \sysname is that it optimizes what the online
agent retains and reasons over, while \sysname aims to bypass online reasoning
for already matched \texttt{Function} segments. Program-guided context
management and cache-first execution are complementary: a host agent can use
program-style memory for misses and long-horizon fallback, while \sysname serves
verified repeated segments as function-cache hits.

\subsection{Verifier-Based and Per-Step Acceleration}

V-Droid changes the role of the model from generator to
verifier~\cite{dai2025advancing}. Instead of asking the LLM to
autoregressively produce an action, it constructs a discretized action space and
asks a trained verifier to score candidate actions. The method includes
prefilling-only verification for speed, pair-wise progress-preference training,
and a scalable human-agent annotation loop. The reported model uses
Llama-3.1-8B and 110K collaboratively annotated samples across four rounds. It
reports 59.5\% success on AndroidWorld, 38.3\% on AndroidLab, and 49\% on
MobileAgentBench; the paper also reports 4.3\,s step-wise latency and 0.74\,s
decision latency in the AndroidWorld comparison table.

V-Droid is strong evidence that mobile GUI agents can be made faster and more
accurate without relying on very large cloud models for every action. Its
limitation for repeated-work reuse is structural: a verifier still evaluates
candidates on each step. A cache-first system asks a different question. For a
matched repeated segment, can the system skip the online decision and execute a
previously accepted \texttt{Function} directly, with only local checks and
effect verification? \sysname should be framed as complementary to V-Droid
rather than as a replacement. A verifier-based agent can serve as the host
fallback for misses, ambiguous matches, unsafe transfers, or novel states.

Other per-step acceleration methods reduce the amount of visual or historical
context sent into the model. ShowUI builds a vision-language-action model for
GUI agents with UI-guided visual-token selection and interleaved
vision-language-action streaming~\cite{lin2025showui}. It reports a 2B model
trained with 256K data, 75.1\% zero-shot screenshot grounding accuracy, 33\%
redundant visual-token reduction during training, and 1.4$\times$ speedup.
FocusUI targets efficient UI grounding by selecting instruction-relevant visual
tokens while preserving positional continuity through
PosPad~\cite{ouyang2026focusui}. It reports that FocusUI-7B improves
ScreenSpot-Pro performance by 3.7\% over GUI-Actor-7B and, with 30\% visual-token
retention, drops only 3.2\% while achieving up to 1.44$\times$ faster inference
and 17\% lower peak GPU memory.

These methods are valuable because high-resolution UI screenshots create heavy
visual-token workloads. Their limitation is that they optimize the current model
call rather than removing the need for the model call on repeated segments. They
help a host agent decide what to do on a screen; they do not by themselves turn
a successful historical run into a callable \texttt{Function} cache entry.
\sysname's benefit should be stated at the segment boundary: a cache hit skips
online reasoning for the matched \texttt{Function}, while per-step accelerators
remain useful for miss handling, fallback, and new or ambiguous states.

\subsection{GUI Grounding Models}

GUI grounding models address the target-localization subproblem. SeeClick
trains a visual GUI agent that relies only on screenshots and introduces
ScreenSpot, a grounding benchmark spanning mobile, desktop, and web
environments~\cite{cheng2024seeclick}. Its key insight is that many GUI agents
depend on long structured representations such as HTML or accessibility trees,
while screenshots are universal across platforms. UGround similarly advocates
visual-only GUI grounding and trains on 10M elements from 1.3M screenshots; it
reports up to 20\% absolute improvement over prior GUI visual grounding models
and shows gains in downstream agent settings~\cite{gou2024uground}. These
models are especially relevant when accessibility metadata is missing, noisy, or
unavailable.

The limitation is runtime and problem scope. A visual grounding model receives
an instruction and current screenshot, then predicts a coordinate or element.
This solves the zero-prior grounding problem. \sysname's transfer setting has a
different prior: the system already has a successful source run, a source
action, neighboring anchors, and a candidate current state. The task is to
rebind a recorded action to the current GUI under a cache-hit contract. A
grounding model could be used as a fallback oracle when structural transfer is
ambiguous, but using a VLM for every cached step would reintroduce the per-step
model cost that cache-first execution is designed to avoid.

This distinction is important for fair related work. GUI grounding papers should
not be described as weak baselines. They solve a broad and difficult
generalization problem from screenshot and language alone. \sysname narrows the
online problem by spending prior execution evidence: a \texttt{Function} carries
source context and effect evidence, \texttt{PageVectorSet} narrows the candidate
state family, and \texttt{control-aware execution} rejects or falls back when the
transfer evidence is insufficient.

\subsection{Benchmarks, Datasets, and What They Measure}

AndroidWorld and AITW support different evaluation questions. AITW measures
learning from large-scale demonstrations across device versions, screen sizes,
applications, and natural-language instructions~\cite{rawles2023android}. It
highlights visual robustness and variation in user commands. AndroidWorld
measures live task execution with programmatic success signals over 116 tasks
and 20 apps~\cite{rawles2025androidworld}. It is stronger for end-to-end
validation because each task has initialization and success-checking logic, but
it still evaluates agents under benchmark-specific task definitions.

MobileViews contributes scale for representation and grounding
analysis~\cite{mobileviews}. Recent descriptions report a million-scale mobile
GUI dataset with over 1.2M unique screens from more than 30K Android apps, and
earlier versions report hundreds of thousands of screenshot-view-hierarchy pairs
over more than 20K apps. Such datasets are useful for studying how often UI
attributes are present, how often WebView or container controls obscure semantic
signals, and how representation methods behave across app families. They do not
directly certify task success unless paired with action, state transition, and
effect labels.

For \sysname, this means the evaluation narrative should keep three measurements
separate. First, representation recall asks whether \texttt{PageVectorSet}
retrieves the correct state family or \texttt{Function} candidates. Second,
action transfer asks whether a recorded action can be safely rebound to the
current GUI. Third, end-to-end success asks whether the whole task reaches a
terminal/effect condition under the benchmark oracle or app-specific verifier. A
high recall number alone is not a live execution guarantee, and a live success
number from one bounded workload is not a production-scale network-effect
result.

\subsection{Positioning of \sysname}

The related work suggests a precise gap. Record-and-replay systems provide
executable reuse, but their memory is tightly bound to the recorded path and
device state. Memory-based GUI agents reuse app knowledge, but much of the reuse
remains prompt or procedure context for online model decisions. Programmatic
agents synthesize executable plans, but they commonly depend on upfront
exploration or coarse task programs that need repair when the live GUI diverges.
Per-step acceleration and visual grounding make each decision cheaper or more
accurate, while repeated workflows still incur online decisions.

\sysname's paper should state a narrow affirmative claim. It treats successful
GUI executions as a cache construction source. A canonical \texttt{run\_log}
becomes a reusable \texttt{Function} asset. A cache lookup retrieves candidate
\texttt{Functions} and compact state evidence. \texttt{PageVectorSet} provides
privacy-bounded state recall without sharing raw screenshots, raw XML, or raw
text as the primary shared object. On a cache hit, Flow checks the start-node
precondition, optionally uses route-only navigation to reach the start state,
transfers recorded actions through source-current evidence, and verifies
terminal/effect evidence under \texttt{control-aware execution}. Misses, unsafe
matches, ambiguous transfers, or failed checks return to the host agent or
fallback path.

This positioning avoids overclaiming. \sysname does not need to be a new
foundation GUI planner. It does not need to make first-run exploration free. It
does not need to prove that \texttt{PageVectorSet} alone guarantees correct
execution. The contribution is a cache-system boundary for repeated GUI work:
accepted execution evidence becomes callable assets, cache hits skip online
reasoning for matched \texttt{Function} segments, and shared reuse is governed
by compact state representation plus verification rather than raw replay.


\begin{table*}[t]
\centering
\small
\caption{Method-level comparison for background and related work. Metrics in
this table are reported by the cited papers and are not \sysname evaluation
results.}
\label{tab:expanded-related-summary}
\begin{tabularx}{\textwidth}{@{}p{0.12\textwidth}p{0.15\textwidth}X X X@{}}
\toprule
\textbf{Direction} & \textbf{Representative work} & \textbf{Method} & \textbf{Reported metrics from source} & \textbf{Limitation for \sysname's target problem} \\
\midrule
Record/replay & RERAN~\cite{gomez2013reran} & Captures and replays low-level Android input and sensor events with timing fidelity. & Focuses on replay fidelity for Android executions. & Tied to recorded timing, device state, and layout; lacks semantic/effect verification for changed GUI state. \\
\addlinespace
Programming by demonstration & SUGILITE~\cite{li2017sugilite} & User demonstrates an Android task; accessibility APIs support later automation from a verbal command. & Demonstrates end-user smartphone automation by demonstration. & Generalization depends on future UI attributes and state compatibility. \\
\addlinespace
App-memory agent & AutoDroid~\cite{wen2024autodroid} & Automated exploration builds app-specific memory and functionality-aware UI representation for LLM execution. & 90.9\% action accuracy and 71.3\% task success on 158 tasks. & Memory guides online decisions; app exploration is a cold-start cost; direct executable cache-hit semantics are outside its main boundary. \\
\addlinespace
Human-like memory & MobileGPT~\cite{lee2024mobilegpt} & Explore, select, derive, and recall; stores modular procedures and app memory. & 82.7\% task learning/automation accuracy; 98.75\% context adaptation; 62.5\% latency and 68.8\% cost reduction vs. GPT-4 baseline. & Recalled memory still needs current-state binding and online adaptation. \\
\addlinespace
Knowledge-base multimodal agent & AppAgent~\cite{zhang2025appagent} & Exploration or demonstrations build a knowledge base for a multimodal smartphone agent. & Tested on 50 tasks in 10 apps. & Knowledge base supports reasoning; it is not a verified \texttt{Function} cache store. \\
\addlinespace
Programmatic state memory & ActionEngine~\cite{zhong2026actionengine} & Crawling Agent builds state-machine memory; Execution Agent synthesizes Python programs. & 95\% success on Reddit/WebArena tasks; about one LLM call on average; 11.8$\times$ cost reduction and 2$\times$ latency reduction. & Depends on offline crawling and program repair; reported setting is web tasks; live mobile heterogeneity and privacy-bounded sharing require separate treatment. \\
\addlinespace
Long-horizon context management & AgentProg~\cite{tian2025agentprog} & Represents interaction history as program variables/control flow with a global belief state. & Reports state-of-the-art results on AndroidWorld and extended long-horizon tasks. & Optimizes retained context for the online agent; repeated matched segments still need a cache-execution boundary to skip online reasoning. \\
\addlinespace
Verifier agent & V-Droid~\cite{dai2025advancing} & Scores discrete candidate actions with a trained verifier and prefilling-only workflow. & 59.5\% AndroidWorld success; 38.3\% AndroidLab; 49\% MobileAgentBench; 4.3\,s step latency; 0.74\,s decision latency in the reported table. & Strong per-step decision method; repeated routines still involve verifier calls at each step. \\
\addlinespace
Visual grounding & SeeClick and UGround~\cite{cheng2024seeclick,gou2024uground} & Screenshot-only GUI grounding to locate actionable elements. & SeeClick introduces ScreenSpot; UGround trains on 10M elements from 1.3M screenshots and reports up to 20\% absolute improvement over prior GUI grounding models. & Solves zero-prior grounding; using it on every cached step reintroduces model inference. \\
\addlinespace
Visual-token acceleration & ShowUI and FocusUI~\cite{lin2025showui,ouyang2026focusui} & Selects or preserves important visual tokens for faster GUI grounding/agent inference. & ShowUI reports 75.1\% zero-shot grounding, 33\% redundant-token reduction, and 1.4$\times$ speedup; FocusUI reports 1.44$\times$ faster inference and 17\% lower peak GPU memory with 30\% token retention. & Reduces model-call cost; cache-first reuse removes model calls for matched \texttt{Function} segments. \\
\bottomrule
\end{tabularx}
\end{table*}


\bibliographystyle{ACM-Reference-Format}
\bibliography{refs}

\end{document}
